% R038：退化情形的全体最优解，以及同一数据上两种回归的不同残差方向。
\begin{center}
\begin{tikzpicture}[x=1cm,y=1cm,>=Stealth,every node/.style={font=\small}]
\path[use as bounding box] (0,0) rectangle (12.4,7.25);
\node[font=\large\bfseries,text=CTrap] at (6.2,6.90) {易错对照：没有横向变化时斜率不唯一；两种回归不能互相倒过来};

% 左：Sxx=0 时的一族最优直线
\draw[rounded corners=5pt,fill=red!2,draw=CTrap] (.15,.40) rectangle (6.02,6.40);
\node[font=\bfseries,text=CTrap] at (3.08,6.02) {$x=x_0$：斜率无法唯一估计};
\draw[->,black!65] (.70,1.35)--(5.42,1.35) node[right] {$x$};
\draw[->,black!65] (.70,1.35)--(.70,4.92) node[above] {$y$};
\draw[red!75!black,densely dashed,thick] (3.00,1.50)--(3.00,4.72);
\foreach \y in {1.85,2.55,3.38,4.25}{\fill[blue!80!black] (3.00,\y) circle (2.8pt);}
\fill[orange!90!black] (3.00,3.00) circle (3.8pt);
\node[anchor=south west,xshift=5pt,text=orange!85!black] at (3.00,3.00) {$(x_0,\bar y)$};
\draw[gray!65,thick] (1.55,1.55)--(4.45,4.35);
\draw[CStatement,very thick] (1.35,3.00)--(4.70,3.00);
\draw[gray!65,thick] (1.55,4.42)--(4.45,1.62);
\node[font=\scriptsize,align=center,text=red!75!black] at (3.00,1.86) {红色竖虚线是所有样本的位置 $x=x_0$\\不是候选回归线};
\node[font=\scriptsize,align=center] at (3.08,.78) {任意 $b$，取 $a=\bar y-bx_0$\\都经过 $(x_0,\bar y)$};

% 右：数据 (1,1),(2,2),(3,2),(4,4)，中心为 (2.5,2.25)
\draw[rounded corners=5pt,fill=red!2,draw=CTrap] (6.40,.40) rectangle (12.25,6.40);
\node[font=\bfseries,text=CTrap,align=center] at (9.32,6.02) {同一数据的两种回归：优化方向不同};
\draw[->,black!65] (6.98,1.35)--(11.72,1.35) node[right] {$x$};
\draw[->,black!65] (6.98,1.35)--(6.98,4.92) node[above] {$y$};
\foreach \x/\y in {7.58/1.75,8.18/2.55,8.78/2.55,9.38/4.15}{\fill[blue!80!black] (\x,\y) circle (2.8pt);}
% y 对 x：hat y=.9x；x 对 y 的等价绘图式 y=19/18 x-7/18。
\draw[CStatement,very thick] (7.58,1.67)--(9.38,3.83);
\draw[orange!85!black,very thick,dashed] (7.58,1.43)--(9.38,3.97);
% 完整的一条纵向残差与一条横向残差，均标明预测点。
\draw[blue!75!black,thick] (8.18,2.55)--(8.18,2.21);
\fill[blue!75!black] (8.18,2.21) circle (2.0pt);
\node[font=\scriptsize,text=blue!75!black,anchor=west] at (8.25,2.18) {$(x_i,\hat y_i)$};
\draw[orange!85!black,thick] (7.58,1.75)--(7.76,1.75);
\fill[orange!85!black] (7.76,1.75) circle (2.0pt);
\node[font=\scriptsize,text=orange!85!black,anchor=west] at (7.83,1.76) {$(\hat x_i,y_i)$};
\fill[orange!90!black] (8.48,2.75) circle (3.8pt);
\draw[orange!85!black,->,thick] (9.13,3.48)--(8.55,2.82);
\node[font=\scriptsize,text=orange!85!black,anchor=west] at (8.55,2.82) {中心};
\node[font=\scriptsize,text=CStatement,align=center] at (9.32,5.35) {$y$ 对 $x$：$\hat y=0.9x$（蓝实线）};
\node[font=\scriptsize,text=orange!85!black,align=center] at (9.32,4.94) {$x$ 对 $y$：$\hat x\approx0.947y+0.368$（橙虚线）};
\node[font=\scriptsize,text=black!75,align=center] at (9.32,.78) {为画在普通 $x$-$y$ 图中，橙线等价写为\\$y=\dfrac{19}{18}x-\dfrac7{18}\approx1.056x-0.389$};
\end{tikzpicture}
\end{center}
